%==============================================================================
%  RELATÓRIO 2  —  Entrega: 13/05/2026
%==============================================================================

\section{Relatório 2: Formulação Matemática e Primeiros Resultados}

\subsection{Introdução}

Este relatório apresenta a segunda etapa do projeto. Com base na concepção
desenvolvida no Relatório 1, são apresentados: a revisão bibliográfica que
fundamenta a técnica escolhida, a formulação matemática completa e validada do
modelo de Programação Linear para o problema de mix de produção da Rentex, e os
primeiros resultados computacionais com análise preliminar.

Os parâmetros do modelo foram refinados em levantamento conduzido junto ao setor
de Planejamento e Controle da Produção (PCP) da Rentex, substituindo as
estimativas conceituais do Relatório 1 por valores reais extraídos do sistema de
gestão da empresa.

\newpage

\subsection{Revisão Bibliográfica}
\label{sec:biblio}

\subsubsection{O Problema de Definição do Mix de Produção}

O problema de definição do mix de produção consiste em determinar as quantidades
a fabricar de cada produto em um horizonte de planejamento fixo, de forma a
maximizar a margem de contribuição total, respeitando os limites de capacidade
dos recursos produtivos e as restrições de demanda de mercado
\cite{hillier2006, arenales2007}.

Do ponto de vista matemático, quando as quantidades são tratadas como variáveis
contínuas, o problema se enquadra na Programação Linear (PL). Quando a
indivisibilidade dos lotes é relevante, trata-se de Programação Linear Inteira
Mista (PLIM) \cite{sobreiro2013}. A PL contínua é adequada para o caso da Rentex:
os produtos são medidos em lotes de 100 metros lineares e, dado o volume semanal
de produção, a hipótese de divisibilidade é válida.

A formulação geral do problema de mix de produção -- apresentada por
\textcite{sobreiro2013} e diretamente adotada neste trabalho -- é:

\begin{align}
  &\text{Maximizar} \quad Z = \sum_{j=1}^{n} (P_j - c_j)\, x_j \nonumber \\[4pt]
  &\text{s.a.} \quad \sum_{j=1}^{n} a_{ij}\, x_j \leq b_i, \quad i = 1, \ldots, m \nonumber \\[2pt]
  &\phantom{\text{s.a.}} \quad x_j \leq d_j^{\max},\quad x_j \geq d_j^{\min}, \quad j = 1, \ldots, n \nonumber
\end{align}

onde $P_j$ é o preço de venda, $c_j$ o custo variável, $a_{ij}$ o consumo do
recurso $i$ por unidade do produto $j$, $b_i$ a capacidade disponível e
$d_j^{\min}, d_j^{\max}$ os limites de demanda do produto $j$. A restrição de
demanda mínima, ausente na formulação clássica mas presente neste modelo, captura
os contratos de fornecimento obrigatório dos forros.

\subsubsection{Programação Linear como Técnica Principal}

A Programação Linear é a técnica dominante para problemas de mix de produção em
horizonte de curto prazo por razões objetivas:

\begin{enumerate}[label=\textbf{(\alph*)}]
  \item \textbf{Garantia de otimalidade global:} o método Simplex encontra o ótimo
        global com certeza, ao contrário de heurísticas que fornecem apenas
        aproximações \cite{hillier2006}.

  \item \textbf{Análise de sensibilidade nativa:} a solução ótima vem acompanhada
        de preços-sombra e intervalos de ranging que permitem avaliar o impacto de
        variações nos parâmetros sem reprocessar o modelo \cite{winston2002}.
        Essa informação é diretamente útil para decisões táticas como negociação
        de horas extras e compra de estoque adicional de fios.

  \item \textbf{Escala e velocidade:} problemas com dezenas de variáveis são
        resolvidos em frações de segundo por solvers modernos \cite{mitchell2011}.
\end{enumerate}

\textcite{sobreiro2013} adotam a solução ótima da PL inteira como referência de
qualidade (\textit{benchmark}) para comparar heurísticas construtivas. Para o
porte do problema da Rentex (4 produtos, 5 recursos), a PL encontra a solução
ótima instantaneamente, tornando desnecessário o uso de qualquer heurística.

\subsubsection{Teoria das Restrições e Identificação de Gargalos}

A Teoria das Restrições (TOC), desenvolvida por \textcite{goldratt2006}, identifica
que o desempenho de qualquer sistema é limitado por sua restrição mais crítica.
No contexto da PL, o gargalo é o recurso com maior preço-sombra na solução ótima:
aquele para o qual uma unidade adicional de capacidade gera o maior incremento na
margem de contribuição total.

A conexão entre TOC e PL é evidenciada pelo trabalho de \textcite{sobreiro2013},
que utiliza o gargalo identificado pela TOC para construir a sequência de
prioridade de produção de sua heurística. Neste projeto, a informação de gargalo
é extraída diretamente da solução ótima da PL, permitindo uma análise gerencial
acessível para a equipe da Rentex.

\subsubsection{Tratamento do Tempo de Setup na Formulação PL}

O Relatório 1 identificou os tempos de setup como um elemento relevante do problema,
especialmente para a Sarja Joy, que requer troca de urdume e trama nos teares.
Na PL, o setup é incorporado por meio do coeficiente técnico ajustado
$a_{1j}^{\text{ajust}}$, definido como:

\begin{equation}
  a_{1j}^{\text{ajust}} = a_{1j}^{\text{ciclo}} + \frac{S_j}{Q_j^{\min}}
  \label{eq:setup}
\end{equation}

onde $a_{1j}^{\text{ciclo}}$ é o tempo de ciclo puro por lote, $S_j$ é o tempo
de setup por troca e $Q_j^{\min}$ é o tamanho mínimo de lote contratual. Essa
abordagem amortiza o setup sobre o lote mínimo, fornecendo um coeficiente
conservador que incorpora o custo do tempo de preparação sem introduzir variáveis
binárias (o que converteria o problema em PLIM) \cite{arenales2007}.

\subsubsection{Justificativa Final da Técnica}

A Programação Linear contínua foi selecionada como técnica principal porque:
(i) a função de margem de contribuição é linear no volume produzido; (ii) as
restrições de capacidade são lineares; (iii) o horizonte semanal valida a hipótese
de parâmetros estacionários; e (iv) a análise de sensibilidade gerada é de valor
gerencial imediato para a Rentex. A Programação Linear Inteira será utilizada
como verificação complementar para confirmar que a hipótese de divisibilidade
não compromete a solução.

\newpage
